\documentclass[../../../include/open-logic-section]{subfiles}

\begin{document}

% Part: first-order-logic
% Chapter: models-theories
% Section: size-of-structures

\olfileid{fol}{mat}{siz}

\olsection{表达 \printtoken{P}{structure} 的大小}

\begin{explain}
有些结构属性，我们甚至无需使用语言的非逻辑符号就能表达。
例如，存在这样一些语句：它们在一个结构中为真，当且仅当该结构的论域至少、至多或恰好包含 $n$ 个元素。
\end{explain}

\begin{prop}
语句
\begin{multline*}
  !A_{\ge n} \ident \lexists[x_1][\lexists[x_2][\dots\lexists[x_n][{}]]]\\
  \begin{aligned}
  (\eq/[x_1][x_2] \land {}
  \eq/[x_1][x_3] \land \eq/[x_1][x_4] \land \dots \land \eq/[x_1][x_n] \land {}\\
\eq/[x_2][x_3] \land \eq/[x_2][x_4] \land \dots \land {} \eq/[x_2][x_n] \land {} \\
\vdots\\
\eq/[x_{n-1}][x_n])
\end{aligned}
\end{multline*}
在结构~$\Struct M$ 中为真，当且仅当 $\Domain M$ 至少包含 $n$ 个元素。因此，$\Sat{M}{\lnot !A_{\ge n+1}}$ 为真当且仅当 $\Domain M$ 至多包含~$n$ 个元素。
\end{prop}

\begin{prop}
语句
\begin{multline*}
!A_{= n} \ident \lexists[x_1][\lexists[x_2][\dots\lexists[x_n][{}]]] \\
\begin{aligned}
  (\eq/[x_1][x_2] \land {}
  \eq/[x_1][x_3] \land \eq/[x_1][x_4] \land \dots \land \eq/[x_1][x_n] \land {}\\
\eq/[x_2][x_3] \land \eq/[x_2][x_4] \land \dots \land {} \eq/[x_2][x_n] \land {} \\
\vdots\\
\eq/[x_{n-1}][x_n] \land {} \\
\lforall[y][(\eq[y][x_1] \lor \dots \lor \eq[y][x_n]]))
\end{aligned}
\end{multline*}
在结构~$\Struct M$ 中为真，当且仅当 $\Domain M$ 恰好包含 $n$ 个元素。
\end{prop}

\begin{prop}
一个结构是无穷的，当且仅当它是
\[
\{!A_{\ge 1}, !A_{\ge 2}, !A_{\ge 3}, \dots \}.
\]
的模型。
\end{prop}

不存在单一的纯粹逻辑语句，使得它在~$\Struct M$ 中为真当且仅当 $\Domain M$ 是无穷的。
然而，我们可以给出一些含有非逻辑谓词符号的语句，它们只有无穷模型（尽管并非每个无穷结构都是它们的模型）。
“是有穷结构”这一性质，以及“是不可数结构”这一性质，甚至无法用无穷的语句集合来表达。
这些事实都是紧致性定理和 L\"owenheim--Skolem 定理的推论。
\end{document}